\uuid{eiIm}
\exo7id{3434}
\titre{exo7 3434}
\auteur{quercia}
\organisation{exo7}
\datecreate{2010-03-10}
\isIndication{false}
\isCorrection{true}
\chapitre{Déterminant, système linéaire}
\sousChapitre{Autre}
\cpge{mpsi-mp2i.calcul-matriciel-systemes.c, pcsi.calcul-matriciel-systemes-lineaires.c, ptsi.calcul-matriciel-systemes-lineaires.c, mpsi-mp2i.matrices-applications-lineaires.b, pcsi.matrices-applications-lineaires.b, ptsi.matrices-applications-lineaires.b, mpsi-mp2i.matrices-applications-lineaires.c, pcsi.matrices-applications-lineaires.d, ptsi.matrices-applications-lineaires.d, bcpst1.systemes-lineaires.a, bcpst1.matrices, bcpst1.applications-lineaires-matrices, bcpst2.applications-lineaires-matrices.a, bcpst2.applications-lineaires-matrices.b, bcpst2.applications-lineaires-matrices.c}
\module{Algèbre}
\niveau{L2}
\difficulte{}

\contenu{
\texte{
On considère un système linéaire homogène : $(S) \iff AX = 0$, avec
$A \in \mathcal{M}_{n,p}(K)$, $n < p$ et $\mathrm{rg} A = n$.
}
\begin{enumerate}
    \item \question{Montrer qu'on peut compléter $A$ en une matrice
     $B = \begin{pmatrix} A \cr A' \cr \end{pmatrix}$ inversible.}
    \item \question{Montrer que les colonnes $n+1$ à $p$ de ${}^{t\!}\text{com}\, B$ constituent
     une base des solutions de $(S)$.}
    \item \question{Consid\'erer l'exemple suivant : 
  $$(S) \iff \left\{
  \begin{array} {lllllllll}
  x &+& 2y &+& 3z &+& 4t &=& 0  \cr
 2x &+& 3y &+& 4z &+& 5t &=& 0. \cr
 \end{array}
 \right.$$}
\reponse{
3. On complète par $\begin{pmatrix} 0 & 0 & 1 & 0 \cr 0 & 0 & 0 & 1 \cr \end{pmatrix}
               \Rightarrow \text{base} =
               \left\{\,
               \left(\begin{smallmatrix} -1 \cr 2 \cr -1 \cr 0 \cr \end{smallmatrix}\right),
               \left(\begin{smallmatrix} -2 \cr 3 \cr 0 \cr -1 \cr \end{smallmatrix}\right)
               \,\right\}$.
}
\end{enumerate}
}
